% ID: SEL-D0039-OB-E01
% ESTADO: validada. S040, cotejada con examen y criterios.
\textbf{a)} Sustituimos $X=2A-Y$ en $X+B=2Y$ y obtenemos
$3Y=2A+B$. Entonces
\[
Y=\frac13\begin{pmatrix}6&3\\3&3\end{pmatrix}
=\boxed{\begin{pmatrix}2&1\\1&1\end{pmatrix}},\qquad
X=2A-Y=\boxed{\begin{pmatrix}2&5\\-3&1\end{pmatrix}}.
\]
La suma da $2A$ y al sumar $B$ a $X$ resulta $2Y$.

\textbf{b)} $A$ es $2\times2$, $C$ es $3\times2$ y
$C^t$ es $2\times3$.
\begin{itemize}
\item $A+D=C$ es imposible porque la suma tendría dimensión $2\times2$.
\item $AD=C^t$ es posible para $D$ de dimensión $2\times3$.
\item $DA=C$ es posible para $D$ de dimensión $3\times2$.
\item $DA=C^t$ es imposible porque $DA$ tiene dos columnas.
\end{itemize}
